\section{Motivation and Challenges}

In this section, we first introduce some key features of mobile NPUs, and then analyze the opportunities of leveraging NPUs' free compute as well as the challenges in implementing efficient system for test-time scaling workloads.

\subsection{Qualcomm's Hexagon NPU} \label{sec:hexagon-npu}

% \ting{this whole section needs to be submilated, extract research-question related features, and leave details to the implementation part.}
% \todo{Add illustrations and tables}

The Hexagon NPU on Qualcomm’s Snapdragon SoC is a representative mobile NPU due to its typical architecture, widespread adoption, and relatively accessible SDK. Therefore, we use it to demonstrate the core features of mobile NPUs.

\subsubsection{Programming Interface}

The primary approach to program Qualcomm's Hexagon NPU is through Qualcomm AI Engine Direct~\cite{qualcomm-qnn} (often referred to as QNN), a proprietary, closed-source DNN inference framework. In most cases, developers cannot customize high-performance low-level kernels even though the full LLVM toolchain for Hexagon NPU is provided in the Hexagon SDK, mainly because the instructions for the matrix unit remain undisclosed. We are able to utilize the FP16 matrix unit by reverse engineering the undocumented instructions in the binary libraries.

\subsubsection{Architecture}

\begin{figure}[h]
    \centering
    \includegraphics[width=1.0\linewidth]{figures/hexagon-npu-arch.pdf}
    \caption{Hexagon NPU Architecture.}
    \label{fig:hexagon-npu-arch}
\end{figure}

\paragraph{Computation Units.} The Hexagon NPU features a typical hybrid architecture of “vector + matrix”. Its vector and matrix units are named HVX (Hexagon Vector eXtension) and HMX (Hexagon Matrix eXtension), respectively. The Hexagon NPU incorporates 6 to 8 scalar VLIW hardware threads for logical control. All vector or matrix instructions are issued from one of the four VLIW slots in a scalar core. The HVX unit context comprises 32 vector registers with a width of 1024 bits, and the number of such units ranges from 4 to 6. The number of HMX units is deduced to be 1 or 2.

\paragraph{Memory Subsystem.} The Hexagon NPU includes a shared 1 MiB L2 cache and 8 MiB of TCM (Tightly Coupled Memory), the latter being a segment of software-managed on-chip memory. The HVX can read data from either the L2 cache or the TCM. Vector scatter/gather operations and all HMX instructions can only access TCM. Data can be loaded from DDR memory into the L2 cache and TCM via the \texttt{l2fetch} instruction and DMA mechanisms, respectively. Both support asynchronous transfers of 1D or 2D tensor data.

% \subsubsection{Computation Units}

% Qualcomm's Hexagon NPU originates from Hexagon cDSP (Compute DSP), which features VLIW ISA and wide SIMD vectors. The Hexagon core is made of several hardware DSP threads, typical 6 or 8 on recent hardware. The DSP hardware threads perform scalar operations and are mainly used for control flows, so we omit them from the following text as they contribute little to the overall computation capabilities. 

% Each DSP hardware thread has access to the Hexagon scalar instructions, which perform fixed-point and floating-point operations on single or pairs of 32-bit registers. These core hardware threads, which we refer to as scalar cores, are omitted from the following text since they are mainly used for control flow and address calculation, contributing little to the overall computation capabilities. (We especially note that the scalar threads are 32-bit cores and that they only support 32-bit virtual addresses. DNN models larger than 4GiB cannot be run directly on Hexagon NPUs without special processing like model chunking.)

% The main arithmetic components on Hexagon NPUs are the HVX (Hexagon Vector eXtension) and HMX (Hexagon Matrix eXtension) units, responsible for vector and matrix computations, respectively. HVX is a coprocessor that adds 128-byte vector processing capabilities to the cDSP. The scalar cores use the HVX coprocessor by accessing an HVX register file, also referred to as the HVX context, which contains 32 1024-bit vector registers and 4 vector predicate registers. Recent DSPs have 4 or 6 HVX contexts. We leave detailed information about the HMX unit to Section~\ref{sec:hmx}.

% \subsubsection{Memory Subsystem}

% The cDSP has a two-level cache subsystem. L1 cache is only accessible to the scalar unit, making L2 the second level cache for the scalar unit and the first level cache for the HVX coprocessor. The L2 cache has a typical size of 1MiB and is shared among the scalar and vector units.

% To hide global DDR memory access latencies, cDSP software usually needs to manually issue data cache fetch/L2 fetch instructions ahead of time. The L2 fetch instruction will not block normal code execution, and it supports prefetching 1D or 2D memory region.

% In addition to regular cache systems, the cDSP also includes a software-managed shared on-chip memory called VTCM (Vector Tightly Coupled Memory). VTCM is a low-power memory that provides approximately twice the bandwidth of L2 and reduces store-to-load latencies. It is also required for certain instructions including vector scatter/gather and all HMX instructions. The typical size of VTCM is 8 MiB, 8 times larger than L2.

% HZX: Remove descriptions about DMA. We don't mention them in this paper.

% \hzx{DMA instructions.} After introducing high-throughput matrix extension instructions, memory accesses through the HVX context may fail to provide sufficient memory bandwidth. Therefore, Qualcomm has simultaneously introduced a set of DMA (Direct Memory Access) instructions. The known DMA instructions involve bidirectional data transfer between the global DDR memory and VTCM with optional L2 cache bypassing functionality, executed asynchronously: software can use the \texttt{dmwait} instruction to block until the DMA engine completes execution, or use \texttt{dmpoll} to poll the completion status of DMA requests in a non-blocking manner. Similarly to the L2 fetch instruction, the DMA instructions also support transferring a 2D memory block.

% \hzx{A table to show bandwidth achieved via HVX vmem load/store and DMA.}

% \subsubsection{The HMX Unit} \label{sec:hmx}

\paragraph{The HMX Unit.} The powerful matrix multiplication capabilities of the Hexagon NPU originate from the HMX component. According to Qualcomm, the HMX unit supports various precisions, including INT4, INT8, INT16, and FP16~\cite{hexagon}. The following introduction is based mainly on FP16 HMX, with relevant information derived from reverse engineering, the Hexagon SDK, the QNN SDK, and publicly available information from Qualcomm.

% However, Qualcomm does not reveal any detailed information about HMX instructions in Hexagon SDK's documents, and the only (official) way to leverage the HMX unit is to use QNN~\cite{qualcomm-qnn}, Qualcomm's proprietary DNN inference framework, which is closed source and mainly targets traditional DNN computation graphs. Fortunately, specific versions of the Hexagon SDK contain an example of FP16 HMX GEMM, whose core HMX functions and layout transformation utilities are provided as binaries. We are able to utilize the FP16 HMX instructions by reverse-engineering these binary libraries. 

The basic data unit for HMX operations is a tile, where each tile contains a matrix of a specific size. For FP16 HMX, a tile measures 32*32, occupying 2 KiB of space. The HMX unit can load several tiles of weight memory and activation memory from the TCM. After performing matrix multiplication on each pair of matrix tiles, it accumulates the results into an internal accumulator. Finally, it outputs a tile corresponding to the accumulator. Meanwhile, the HMX unit can independently scale and add biases to each channel (column) of the output tile.

FP16 HMX tiles have a special memory layout, as shown in Figure~\ref{fig:hmx-fp16-crouton-tile-layout} (a). Both input and output tiles follow this layout. A typical way to construct this layout is to use HVX instructions to perform cross-lane shuffling on every two adjacent rows of the original matrix.

\begin{figure}
    \centering
    \includegraphics[width=0.9\linewidth]{figures/hmx-fp16-crouton-tile-layout-2.pdf}
    \caption{(a) The memory layout of FP16 HMX tile. Each tile corresponds to a 32 * 32 matrix and takes up 2048 bytes. Every two rows are permuted, having the same layout as the transposed 2 * 32 sub-matrix. (b) The overall memory layout for HMX-based GEMM. The weight tiles are arranged in column-major layout since the hardware performs inner-product at tile level.}
    \label{fig:hmx-fp16-crouton-tile-layout}
\end{figure}

% The basic data unit for HMX operations is a tile, where each tile contains a matrix of a specific size. For FP16 HMX, a tile measures 32*32, occupying 2 KiB of space. The tile sizes for other data types vary, but the available information suggests that they all occupy 2 KiB. The HMX unit can load several tiles of weight memory and activation memory from the VTCM. After performing matrix multiplication on each pair of matrix tiles, it accumulates the results into an internal accumulator. Finally, it outputs a tile corresponding to the accumulator. Meanwhile, the HMX unit can independently scale and add biases to each channel (column) of the output tile. \hzx{Possibly related to quantization granularity and precision issue} The instruction names indicate that there might be a certain type of conversion function during the process of outputting data from the internal accumulator. It is important to note that the HMX does not introduce additional registers, and all input and output tiles, as well as scaling and bias data, must be located in the VTCM. The internal accumulator of the HMX has a wider bit-width than the inputs and outputs, but the high-precision intermediate accumulation results cannot be directly fetched to the register file or the VTCM. \hzx{related to output precision issue}

% HZX: too many details. Maybe these information should appear at appendix
% At the instruction level, HMX instructions can be mainly divided into input tile loading instructions, output tile instructions, scaling and bias factor loading instructions, accumulator conversion instructions, and others. Due to our limited understanding of the accumulator conversion instructions, we will skip their description. Input tile loading instructions consist of activation loading and weight loading instructions, which must be paired together to appear in the same VLIW instruction packet. There are numerous instruction variants with different modifiers for each, but they all use two general-purpose registers to indicate the base address and length of the data tile in the VTCM. Once input tile loading instructions are issued, the HMX unit automatically reads each pair of tile data, performs matrix multiplication, and accumulates the results into the accumulator. Our experiments show that a single specific FP16 tile loading instruction can operate on up to 32 contiguous tiles in the VTCM. However, errors occur when these contiguous tiles cross the 1 MiB boundary of the VTCM. Multiple sets of tile-loading instructions can be executed consecutively, and the corresponding matrix products will continue to accumulate until an output tile instruction appears. This type of instruction outputs the internal state to the VTCM and clears the internal accumulator. The scaling and bias factor loading instructions require the base address of a 256-byte area in the VTCM, which contains scaling factors and biases that independently control each output channel.

% Since the HMX unit maintains the internal state of the accumulator across instructions, the Hexagon SDK requires that only one thread can execute HMX instructions at a time. The software must strictly ensure that the HMX unit is shared by different threads in a time-sharing manner. From a software perspective, the HMX hardware unit is globally unique (although there may be multiple physical HMX units).

% \begin{figure}
%     \centering
%     \includegraphics[width=0.8\linewidth]{figures/hmx-instructions-workflow.png}
%     \caption{A sequence of HMX instructions performing matrix tile multiplication, accumulation and output tile channel scaling. Unlike common GEMM-acceleration ISAs, HMX inputs and outputs reside in VTCM, and the accumulator content is not accessible from registers.}
%     \label{fig:hmx-instructions-workflow}
% \end{figure}

\subsection{Opportunities: Free Matrix Computation During LLM Decoding}

During the autoregressive generation process, LLM's input typically corresponds to only one token, which causes the GEMM operation to degenerate into GEMV. For example, an activation matrix of shape [1, \texttt{hidden\_dim}] is multiplied by a weight matrix of shape [\texttt{hidden\_dim}, \texttt{proj\_dim}]. In the case of using FP16 HMX, the effective size of each compute tile is [1,32]$\times$[32,32]. Since the basic unit of hardware computation is a $32\times32$ tile, 31 rows in the input activation tile do not correspond to actually useful content, resulting in low utilization of the matrix unit and waste of computing power.

Meanwhile, some test-time scaling algorithms can achieve better generation quality by increasing the computation during generation, including parallel sampling methods such as Self-Consistency, Best-of-N and Beam Search. Their characteristic is that they explore multiple generation paths using a batch size greater than 1 and use certain ways (e.g., an external verifier) to select the better generation paths. Figure~\ref{fig:tts-example} shows an example of test-time scaling using Best-of-N. As the generation budget (i.e. the maximum batch size in the decoding phase) increases, the model accuracy on the MATH500 dataset improves significantly.

Based on these, we propose running the test-time scaling workloads of LLM on mobile NPUs. In this way, the computing power of the NPU wasted during the conventional LLM generation process can be effectively leveraged. The decoding overhead will not increase significantly in theory, and the generation quality of the model can be improved at run-time without modifying the model weights.

\begin{figure}[h]
    \centering
    \includegraphics[width=0.8\linewidth]{figures/tts-example.pdf}
    \caption{An example of test-time scaling with two models. The accuracy on MATH500 improves as generation budget increases.}
    \label{fig:tts-example}
\end{figure}

% \ting{this section is important, needs to be elevated.}
% It is known that during the decoding phase of LLM inference, GEMV operations dominate the computation, making the matrix multiplication units on the Hexagon NPU underutilized. 

% Due to the significant performance gap between the HMX and the HVX on the Hexagon NPU, it is impossible to achieve efficient GEMV using only the non-HMX part \ting{what is the meaning of this sentence?}. Therefore, conventional implementations also use the HMX unit for GEMV, leading to a waste of computing power. \ting{here needs quantitative data to show. We need 1) underutilized NPU HMX data and 2) test-time scaling data to show the benefit of enlarging batch size. }

% For this "free" computing power provided by the HMX, we can leverage certain test-time scaling methods to enhance the model's inference performance. Specifically, such test-time scaling algorithms include PRM-guided beam search, Best-of-N sampling, and other repeated sampling methods. Their characteristic is that they explore multiple paths using a batch size greater than 1 during generation and use an external verifier to select the better generation paths. We believe that the overhead of calling the external verifier is much smaller compared to the generation itself (e.g., PRM can be efficiently executed with prefix-cached prefill) \todo{provide some data}, and moderately increasing the batch size on the NPU does not significantly increase the generation cost.

\subsection{Challenges}

Although it is theoretically feasible to utilize mobile NPUs for test-time scaling, an efficient implementation faces numerous hardware challenges. We summarize these challenges as follows.

\paragraph{Insufficient Precision.} Although HMX units support FP16 GEMM, deploying FP16 models on resource-constrained devices remains impractical, making quantized models the typical alternative. The matrix units in most mobile NPUs, including the HMX, were originally designed to accelerate integer-quantized DNN models that employ coarse-grained quantization schemes, such as per-tensor or per-channel quantization. As a concrete example, Hexagon NPUs lack native hardware support for the fine-grained quantization methods essential to modern LLMs. This limitation is further reflected in the software stack: QNN only supports per-tensor or per-channel weight quantization. Applying coarse-grained low-bit quantization directly to LLM weights can lead to significant accuracy degradation.

As shown in Table~\ref{tab:llama3.2-1b-performance-compare}, the accuracy results of the Llama 3.2 1B-Instruct model under QNN’s per-channel quantization\footnote{We use the official model released by PowerServe, available at https://huggingface.co/PowerServe/Llama-3.2-1B-PowerServe-QNN29-8G3} and AWQ per-group 4-bit quantization (both under W4A16 settings) indicate that the per-channel quantized model suffers severe performance degradation in challenging mathematical reasoning tasks. Unfortunately, since test-time scaling methods are applied in such tasks, the baseline accuracy achieved by QNN fails to meet even the minimal requirements for performance scaling.

\begin{table}[h]
    \centering
    \begin{tabular}{cccc}
        \hline
        dataset & AutoAWQ (W4A16) & QNN (W4A16) \\
        \hline
        MATH500 ($\uparrow$) & 15.9 & 2.1 \\
        GSM8K ($\uparrow$) & 32.6 & 3.4 \\
        Wiki PPL ($\downarrow$) & 19.42 & 28.99 \\
        \hline
    \end{tabular}
    \caption{Comparison of Llama3.2-1B-Instruct's performance under different implementation. QNN's quantization drastically hurts model's reasoning ability.}
    \label{tab:llama3.2-1b-performance-compare}
    \vspace{-1em}
\end{table}

\paragraph{Weak General Purpose Compute and Memory Bandwidth.} In the absence of native hardware support for fine-grained group quantization, a common approach is to rely on general-purpose computing units to handle such computations. However, we discover a significant gap between the compute and memory access capabilities of the general-purpose vector units and the specialized matrix units within the NPU. We measure the FP16 GEMM performance of both the HVX and the HMX on the Hexagon V75 NPU using a 1024$\times$1024$\times$1024 GEMM operation, with all inputs and outputs residing in on-chip TCM to reflect the hardware's peak performance. As shown in Table~\ref{tab:hexagon-npu-profile}, the FP16 GEMM throughput of the matrix unit reaches up to 12 TFLOPS — over 300 times higher than that of a single vector thread. In terms of memory bandwidth, a dedicated DMA engine achieves over 60 GB/s read bandwidth from DDR, whereas the vector unit's memory read bandwidth via the core data path remains below 30 GB/s. However, the high bandwidth provided by DMA is restricted to large, regular 1D or 2D data blocks and cannot efficiently handle small or irregular memory accesses. These observations highlight that the general-purpose compute and memory bandwidth of the vector unit are insufficient to keep up with the computational throughput of the specialized matrix unit, posing a major challenge for implementing high-performance mixed-precision GEMM kernels under fine-grained quantization.

\begin{table}[h]
    \centering
    \begin{tabular}{ccc}
        \hline
        hardware units & HVX (1 Thread) & HMX \\
        \hline
        FP16 GEMM GFLOPs & 32.93 & 12032.54 \\
        memory read bw. (GB/s) & 26 & 60 (DMA) \\
        \hline
    \end{tabular}
    \caption{The performance metrics of the HVX and HMX units, in terms of FP16 GEMM computing power (in GFLOPs) and memory read bandwidth (in GB/s).}
    \label{tab:hexagon-npu-profile}
    \vspace{-1em}
\end{table}

% \ting{this part should be sublimated too. just keep 2-3 challenges.}
% However, implementing real test-time scaling on the Hexagon NPU faces numerous challenges. TTS algorithms are mainly designed for generative reasoning tasks, which have high demands for model precision. \ting{what is the meaning of last sentence.} Unfortunately, hardware designs such as the HMX have limitations in precision, and the quantized models provided by Qualcomm officially cannot meet usable standards at all. To address this, we adopt fine-grained group quantization for weights and utilize the FP16 HMX to improve the model precision. However, there is a mismatch between the small group size and the large width of the HVX registers,\ting{this can be one challenge} which poses difficulties in implementing an efficient dequantization kernel. Moreover, the softmax operation in the FP16 HMX-centric Attention has both precision and performance issues: the input is forcibly truncated to FP16, and the exponential calculation using polynomial approximation on the HVX is extremely inefficient.

% \todo{Revise and restructure the draft,} We identify several challenges that hinder running LLM workloads on existing commodity mobile NPUs: 

% \begin{itemize}[leftmargin=*, itemindent=1pt]
% \item{Weak General Purpose compute and memory bandwidth}

% The scalar units and vector units of Hexagon NPU cannot provide enough compute and memory bandwidth to match the throughput of the matrix units. On NVIDIA GPUs, implementing memory-bound kernels like GEMV via programming CUDA cores is often more efficient than using TensorCores. Recent works (like SageAttention) targeting NVIDIA GPUs are able to leverage CUDA cores to perform fine-grained operations with negligible overheads. However, this is not the case on the NPU. \ting{this part is not clear. What is mismatch?}

% \item{Granularity mismatch}

% \sout{1. Recent matrix multiplication accelerator hardware usually prefers large block/tile sizes. However, during the decoding phase of LLM generation, GEMV dominates the computation, making the GEMM hardware units underutilized.} \ting{I think you need to specify the LLM inference process first to explain the challenges. why here only mentions SIMD?}

% 2. Quantization algorithms prefer small group sizes to minimize accuracy loss, while Hexagon NPU's SIMD unit operates on large vectors. Processing small groups with padding data in vector registers further reduces computational throughput. Compared to GPU's SIMT architecture, wide vector SIMD is much less flexible.

% \item{Precision gap}

% The HMX unit is designed for accelerating quantized/low-precision operations such as convolutions and matrix multiplications, supporting precisions including INT4, INT8, INT16, and FP16. Unfortunately, directly adopting coarse-grained low-precision operations causes significant damage to the inference performance of LLMs, and the hardware does not exhibit fine-grained compensatory designs for precision. We only observe that the HMX unit can perform independent scaling and bias addition on each column of the output tile, while Qualcomm's QNN framework only supports weight quantization at the per-tensor or per-channel granularity. Table~\ref{tab:llama3.2-1b-performance-compare} shows the accuracy results of Llama 3.2 1B-Instruct with QNN per-channel and AWQ per-group 4-bit quantization, revealing that the per-channel quantized model suffers severe performance degradation in more challenging reasoning tasks. We speculate that the HMX unit was initially designed for CNN workloads, which are less sensitive to precision.

% HZX: We remove some designs so the output precision issue is not relevant now

% Although using FP16 HMX can avoid some issues caused by integer-based HMX, FP16 HMX still has precision limitations: the floating-point intermediate accumulation results of HMX cannot be retrieved with 32-bit precision, and only FP16 tile results can be output ultimately. In contrast, other modern matrix multiplication extension units (such as NVIDIA TensorCore, Intel AMX, and Huawei Ascend Cube unit) almost all allow the accumulation results of FP16/BF16 matrix multiplications to be output to registers with FP32 precision. Based on some measurements, we speculate that FP16 HMX has an FP32-precision accumulator, meaning its output is truncated to FP16. This characteristic of HMX poses challenges for implementing precision-sensitive operators like FlashAttention, which require tile-wise computation.


% \begin{table}
%     \centering
%     \begin{tabular}{cccc}
%         \hline
%         dataset & torch (BF16) & AutoAWQ (W4A16) & QNN (W4A16) \\
%         \hline
%         MATH500 & 24.5 & 15.9 & 2.1 \\
%         GSM8K & 42.5 & 32.6 & 3.4 \\
%         Wiki PPL & 17.80? & 19.42? & 28.99? \\
%         \hline
%     \end{tabular}
%     \caption{Comparison of Llama3.2-1B-Instruct's performance under different implementation. QNN's quantization drastically hurts model's reasoning ability.\ting{don't show FP16 I think?}}
%     \label{tab:llama3.2-1b-performance-compare}
%     \vspace{-1em}
% \end{table}

% \end{itemize}
